\documentclass[10pt,a4paper]{amsart}
\usepackage[margin=19mm]{geometry}
\usepackage{amsmath,amssymb,mathtools}
\usepackage[scr=rsfs,cal=euler]{mathalfa}
\usepackage{xfrac}
\usepackage[unicode,hidelinks,bookmarks=false]{hyperref}
\newcommand{\ie}{i.e.\ }
\newcommand{\cf}{cf.\ }
\newcommand{\resp}{respectively\ }
\newcommand{\Subsection}{\subsection}
\providecommand{\nobreakdash}{\nobreak\hbox{-}}
\setlength{\emergencystretch}{2em}
\multlinegap0pt
\usepackage{dsfont}
\usepackage{bm}
\usepackage{cancel}
\usepackage{amssymb}
\usepackage{mathtools}
\usepackage{tikz}
\usepackage{float}
\usepackage[shortlabels]{enumitem}
\usepackage{algorithm}\usepackage{algpseudocode}
\usepackage{xcolor}
\newtheorem{proposition}{Proposition}
\usepackage{cleveref}
\crefname{proposition}{Proposition}{Propositions}
\DeclareMathOperator{\E}{\mathbb{E}}
\DeclareMathOperator{\Q}{\mathbb{Q}}
\DeclareMathOperator{\R}{\mathbb{R}}
\DeclareMathOperator{\calB}{\mathcal{B}}
\DeclareMathOperator{\calD}{\mathcal{D}}
\DeclareMathOperator{\calF}{\mathcal{F}}
\DeclareMathOperator{\calO}{\mathcal{O}}
\title{\vspace{-10mm} Occupied Processes: Going with the Flow \vspace{-3mm}}
\author{Valentin Tissot-Daguette}
\date{Source: arXiv:2311.07936v6 (CC BY 4.0)}
\begin{document}
\maketitle

\noindent\emph{Source and licence.} \vspace{-10mm} Occupied Processes: Going with the Flow \vspace{-3mm}. Version arXiv:2311.07936v6 (CC BY 4.0), distributed by arXiv under the Creative Commons Attribution 4.0 International licence, http://creativecommons.org/licenses/by/4.0/, as recorded in the arXiv licence metadata for this exact version and shown on the abstract page https://arxiv.org/abs/2311.07936v6. Full licence notice: \texttt{licenses/arxiv-2311.07936-CC-BY-4.0.txt}.\par\smallskip

\noindent\textit{Editorial scope.} Counted unit: the article's proposition prop:Markov (source0.tex lines 540--542). The article's complete original proof of it is delivered with it (source0.tex lines 543--553, one \texttt{proof} environment of 11 lines). No mathematics is written, repaired or re-derived: the statement and the proof are the article's own lines, in the article's own notation, and no retained line is substituted or re-flowed.  Retained uncounted as the source-local context the proof needs: lines 444--462.  Nothing was omitted from any retained span, so the package carries no omission record.  No retained line contains a citation, so the wrapper carries no bibliography and no bibliography entry.  No retained line contains a figure environment, an image-inclusion command or drawing code, so no image is delivered and the package declares no asset dependency.  Editorial formatting only, in addition: an AMS \texttt{amsart} preamble that restates the article's own declarations and macros used by the retained lines, namely the environments \texttt{proposition} on separate counters, together with the standard packages those lines need.  The retained environments print in this document as Proposition 1 in order of appearance, whereas the article prints the same items with the numbering of its own full document.  The complete native source of the article as distributed in the arXiv e-print for this exact version is bundled unchanged as \texttt{sources/149963/source0.tex}.
\begin{proposition}\label{prop:OccProperties} 
The occupation flow $\calO$ satisfies the following properties. 
    \begin{itemize}
        \item[\textnormal{I.}] \textnormal{(Strong Time Additivity)} 
        For all stopping times $\varsigma \le \tau$ and $A \in \calB(\R)$,  
\begin{equation}\label{eq:additive}
    \calO_{\tau}(A) = \calO_{\varsigma}(A) + \calO_{\varsigma,\tau}(A), \quad \Q-a.s., 
\end{equation}
with the incremental flow
$\calO_{\varsigma,\tau}(A) := \int_\varsigma^\tau \mathds{1}_{A}(X_s)d\langle X \rangle_s$. 

        \item[\textnormal{II.}] \textnormal{(Absolute Continuity)} $\calO_t \ll \lambda$ $ \ \Q-$a.s., where $\lambda$ is the Lebesgue measure on $\R$. 
             \item[\textnormal{III.}] \textnormal{(Occupation Time Formula)} For every Borel function $\phi:\R\to \R_{+}$, 
\begin{equation}\label{eq:OTF}
    \int_0^t \phi(X_s)d\langle X \rangle_s = \int_{\R} \phi(x)\calO_t(dx).  
\end{equation}
             
    \end{itemize}
\end{proposition}

\begin{proposition}\label{prop:Markov}
Let $X$ be a strong Markov process. Then the occupied processes $(\calO,X)$,  $(\tilde{\calO},X)$ are  strong Markov as well. 
\end{proposition}

\begin{proof}
This follows from 
the strong time additivity  of the occupation measure (\cref{prop:OccProperties} I): for  
all stopping times $0\le \tau\le \varsigma \le T$ and occupation functional   $f:\calD \to \R$, \begin{align*}
    \E^{\Q}[f(\calO_{\varsigma},X_{\varsigma}) \ |  \calF_{\tau}]
    =  
    \E^{\Q}[f(\calO_{\tau}+\calO_{\tau,\varsigma},X_{\varsigma}) \ |   \calF_{\tau}]
     = \E^{\Q}[f(\calO_{\varsigma},X_{\varsigma}) \ |  (\calO_\tau,X_\tau)], 
\end{align*} 
noting that  $\calO_{\tau,\varsigma}$ (and $X_{\varsigma}$) is conditionally independent of $\calF_{\tau}$ given $X_{\tau}$. The same arguments apply to the calendar time occupied process. 
\end{proof}

\end{document}
